% !TeX program = lualatex
% Compile from this directory: latexmk -norc -lualatex -outdir=paper_build paper.tex
\documentclass[11pt,a4paper]{article}
\usepackage[T1]{fontenc}
\usepackage{lmodern}
\usepackage[margin=28mm]{geometry}
\usepackage{amsmath,amssymb,amsthm,mathrsfs,slashed}
\usepackage{booktabs,array}
\usepackage{microtype}
\usepackage[numbers,sort&compress]{natbib}
\usepackage{xcolor}
\definecolor{linkblue}{RGB}{25,65,115}
\usepackage[colorlinks=true,allcolors=linkblue,breaklinks=true]{hyperref}
\hypersetup{pdftitle={Warped AdS3 backgrounds in eleven-dimensional supergravity: an attributed reconstruction},pdfauthor={}}
\newcommand{\dd}{\mathrm{d}}
\newcommand{\AdS}{\mathrm{AdS}}
\newtheorem{proposition}{Proposition}[section]
\newtheorem{theorem}[proposition]{Theorem}
\numberwithin{equation}{section}
\allowdisplaybreaks[2]
\setlength{\emergencystretch}{1.5em}
\title{Warped $\AdS_3$ backgrounds in eleven-dimensional supergravity:\\[3pt]
\large An attributed reconstruction of a three-fibre sector}
\author{}
\date{}
\begin{document}
\maketitle
\vspace{-2em}
\renewcommand{\abstractname}{Reproduced source abstract}
\begin{abstract}
we construct new families of supersymmetric solutions of eleven-dimensional supergravity by deforming the half-BPS $\displaystyle{\mathrm{AdS}_{3}\times S^{3}_{+}\times S^{3}_{-}\times \mathbb{R}^{2}}$ background. we consider lightlike, spacelike and timelike warpings of the anti-de Sitter factor, accompanied by deformations of the two three-spheres, and determine the supporting four-form fluxes and the constraints imposed by the equation of motion. in particular, the deformation parameters of the two spheres are necessarily equal, although their radii need not coincide. we compute the conserved charges and determine the preserved supersymmetries and Killing superalgebras. generic lightlike solutions preserve four real supercharges, whereas spacelike and timelike solutions preserve eight, with a Killing superalgebra containing a single exceptional $\displaystyle{\mathfrak{d}(2,1;\alpha)}$ factor.
\end{abstract}
\begin{quote}\small
\textbf{Provenance.} The supplied abstract above is retained verbatim, including its transcription conventions. It reproduces the abstract of Noris, arXiv:2611.21727v1~\cite{Noris2026}. The present text is an attributed reconstruction of a specified sector, with complementary calculations; it is not a claim of independent discovery of the backgrounds. The precise ansatz, chirality choices and meaning of ``charges'' are fixed below.
\end{quote}

\section{Introduction}
Warped anti-de Sitter backgrounds provide a controlled setting in which an anisotropic geometry, its supporting flux, and its supersymmetry can be studied together. In eleven dimensions, the product $\AdS_3\times S^3_+\times S^3_-\times\mathbb R^2$ is particularly useful: its two independent sphere radii enter a continuous exceptional-superalgebra parameter, while the field equations impose one relation among the three curvature scales. Its brane interpretation and enhanced near-horizon symmetry are established features of intersecting-brane geometries~\cite{BoonstraPeetersSkenderis1998,GauntlettMyersTownsend1999}; the dependence of the exceptional algebra on the radius ratio is treated explicitly in Ref.~\cite{BachasEtAl2014}.

Deformations generated by commuting isometries supply a general route to flux-supported anisotropy~\cite{LM2005,GauntlettLeeMateosWaldram2005,CatalOzerDeger2009}. Warped and Schr\"odinger geometries also arise in lower-dimensional string-theory constructions~\cite{OrlandoUruchurtu2010,BobevVanRees2011}. These precedents motivate a direct eleven-dimensional calculation: a transformation rule alone does not determine whether all Killing spinors have been found, whether a timelike deformation is globally causal, or which flux integrals define quantized charges.

The immediate source of this reconstruction is Noris~\cite{Noris2026}. Its non-null families have independent AdS and common-sphere deformation parameters. We reconstruct their equal-fibre boundary, rather than the full two-parameter families, and a round-sphere null sector. The resulting exposition derives the field equations, exact local supersymmetry and normalized odd brackets in one convention. The opposite null Hopf choices and an unequal-sphere example delimit the conclusions; they are not presented as a classification of the source's broader ansatz.

Three distinctions organize the calculation. First, equality of the squashing parameters is a consequence of a specified six-component flux ansatz, not a universal requirement of supersymmetry. Second, a projector supplies a candidate initial-value space; necessity follows from supercurvature, and sufficiency requires an invariant flat subbundle of the full superconnection. Third, local flux densities, compact-cycle Page charges, and asymptotic gravitational charges are different observables. We compute the first two, imposing an additional torus quotient only when it is actually used. No warped energy, angular momentum or central charge is inferred without boundary data.

The non-null family is developed in Section~\ref{sec:geometry}, with its exact supersymmetry and exceptional algebra in Section~\ref{sec:susy}. Section~\ref{sec:null} treats the null sector without a singular fibre normalization. Section~\ref{sec:charges} derives the flux charges and global restrictions. The explicit unequal-sphere example in Appendix~\ref{app:unequal} shows why the ansatz qualification is essential.

% !TeX root = paper.tex
\section{Homogeneous deformations and the flux constraint}\label{sec:geometry}
\subsection{Theory, orientation and seed geometry}
We work in the classical, source-free, two-derivative theory~\cite{CremmerJuliaScherk1978}, with mostly-plus signature and orientation $e^{0\cdots10}$. For an orthonormal coframe, $e^{i_1\cdots i_p}$ denotes its ordered wedge product. Our equations are
\begin{align}
\dd F&=0,& \dd\star F+\tfrac12F\wedge F&=0,\notag\\
R_{MN}&=\frac1{12}\left(F_{MPQR}F_N{}^{PQR}-\frac1{12}g_{MN}F^2\right),&
F^2&=F_{MNPQ}F^{MNPQ}.\label{eq:eom}
\end{align}
The Killing-spinor equation is
\begin{equation}
\mathscr D_M\epsilon=\left(\nabla_M+\frac1{12}\Gamma_M\slashed F
-\frac14\slashed{\iota_M F}\right)\epsilon=0,
\qquad \slashed F=\frac1{4!}F_{NPQR}\Gamma^{NPQR}.
\label{eq:kse}
\end{equation}
Equivalently, its flux term is $(\Gamma_M{}^{NPQR}-8\delta_M^N\Gamma^{PQR})F_{NPQR}/288$, with the five-index gamma matrix \emph{fully antisymmetrized}. We take $\Gamma^{0\cdots10}=-1$. The orientation, this central Clifford sign, and the sign of the Chern--Simons term are held fixed together. Appendix~\ref{app:clifford} gives an explicit representation and Majorana structure.

Indices are grouped as $(012)_A,(345)_+,(678)_-,9=x,10=y$. The seed is
\begin{align}
g_0&=g_A(L)+g_+(R_+)+g_-(R_-)+\dd x^2+\dd y^2,\notag\\
F_0&=(f_Av_A+f_+v_++f_-v_-)\wedge\dd x,
& f_A&=\frac2L,\qquad f_\pm=\frac2{R_\pm},\label{eq:seed}\\
L^{-2}&=R_+^{-2}+R_-^{-2}.\label{eq:radii}
\end{align}
All radii are finite and positive. The $v_I$ are physical volume forms of the \emph{seed}, so that $\int_{S^3_\pm}v_\pm=2\pi^2R_\pm^3$. In particular, equality of radii is not assumed. The flat Einstein equations give $F_0^2=24(-f_A^2+f_+^2+f_-^2)=0$; the curved factors then have Ricci eigenvalues $-2/L^2,2/R_+^2,2/R_-^2$. The other equations follow from closure of the factor volume forms and $F_0\wedge F_0=0$.

For the spinorial calculations we use
\begin{gather}
a=L^{-1},\qquad b=R_+^{-1},\qquad c=R_-^{-1},\qquad a^2=b^2+c^2,\notag\\
A=\Gamma^{012},\quad B=\Gamma^{345},\quad C=\Gamma^{678},\quad
X=\Gamma^9,\quad Y=\Gamma^{10},\notag\\
H=aA+bB+cC,\qquad Q=\frac baAB+\frac caAC,\qquad
P_H=\frac{1-Q}{2},\qquad \Pi_0=P_H\frac{1+X}{2}.\label{eq:clifford-data}
\end{gather}
Here $A^2=1$, $B^2=C^2=-1$ and the three volume elements anticommute. Consequently $H^2=0$, $Q^2=1$, and $H\epsilon=0$ is equivalent to $Q\epsilon=-\epsilon$.

\subsection{The non-null three-fibre family}
Choose orthonormal Maurer--Cartan forms on each seed factor, with
\begin{equation}
\dd\theta^i=-\frac12C^{(0)i}{}_{jk}\theta^{jk},\qquad
C^{(0)i}{}_{jk}=f_I\eta^{ii}\varepsilon_{ijk}
\quad(i,j,k\text{ in factor }I).
\label{eq:mc}
\end{equation}
The ordered triples defining $\varepsilon$ are $(012),(345),(678)$. Select sphere fibres $5,8$, and AdS fibre $f=2$ for spacelike warping or $f=0$ for timelike warping. Set $\delta=\eta_{ff}=\pm1$. In particular, the timelike convention has $\dd\theta^0=+f_A\theta^{12}$ while $\dd\theta^5=-f_+\theta^{34}$ and $\dd\theta^8=-f_-\theta^{67}$. These relative choices are part of the supersymmetric ansatz.

Let $g_{\rm base}$ denote the seed metric with these three fibre squares removed, including both flat directions, and put $\Theta=\theta^f\wedge\theta^5\wedge\theta^8$. For real $\widehat\gamma$ define
\begin{align}
G&=(1+\delta\widehat\gamma^2)^{-1}>0,\notag\\
g'&=G^{-1/3}g_{\rm base}
+G^{2/3}\big[\delta(\theta^f)^2+(\theta^5)^2+(\theta^8)^2\big],\notag\\
F'&=F_0+\delta\widehat\gamma G\,\dd\Theta.\label{eq:non-null}
\end{align}
The three fibre-to-base length ratios coincide. It is useful to remove a constant overall scale:
\begin{gather}
s=G^{-1/6},\qquad \lambda=\sqrt G,\qquad
\nu=\widehat\gamma\sqrt G,\qquad \lambda^2+\delta\nu^2=1,\notag\\
\widetilde g=s^{-2}g'=g_{\rm base}+\lambda^2
\big[\delta(\theta^f)^2+(\theta^5)^2+(\theta^8)^2\big],\qquad
\widetilde F=s^{-3}F'=\lambda F_0+\delta\lambda^2\nu\,\dd\Theta.
\label{eq:scaled}
\end{gather}
This is the classical scaling symmetry $(g,F)\mapsto(s^2g,s^3F)$, not a metric rescaling at fixed flux. The orthonormal frame of $\widetilde g$ is $e^i=d_i\theta^i$, where $d_i=\lambda$ on the three fibres and $d_i=1$ otherwise. In the spacelike branch $0<\lambda\leq1$ and $|\nu|<1$; in the timelike branch $\lambda\geq1$ and $\nu\in\mathbb R$, equivalently $|\widehat\gamma|<1$. The endpoint $\lambda=0$ is excluded.

For comparison with the source, the scale-removed metric has $\lambda_{\rm N}=\omega_{\rm N}=\lambda$ in Eqs.~(3.13) and (3.21) of Ref.~\cite{Noris2026}. With $R_{0,\rm N}=L$, $A_0=2\theta^f$ and $A_\pm=2\theta^{5,8}$, its boundary flux is obtained\footnote{Take $x=z_1$, $y=-z_2$ and $v_I=-R_I^3v^{(\rm N)}_{I,0}$, where $v^{(\rm N)}_{I,0}$ is the source's unit-radius undeformed volume form. Reversing one horizontal leg in each three-factor, together with $y$, preserves the eleven-dimensional orientation and gives \eqref{eq:mc}. The electric-leg ordering contributes the compensating minus sign.} by taking $\nu=\sqrt{1-\lambda^2}$ in the spacelike case and $\nu=-\sqrt{\lambda^2-1}$ in the timelike case. The more general terms proportional to $\sqrt{\lambda_{\rm N}^2-\omega_{\rm N}^2}$ are not reconstructed here.

\subsection{Existence and a restricted necessity theorem}
The field equations can be reduced while initially allowing independent squashing ratios. This makes explicit which equality is forced, rather than built into the starting geometry. Let $h_I=q_I\theta_I$ be the three orthonormal fibre forms, $q_I>0$, and $\Omega_I$ their horizontal area forms. Choose
\begin{equation}
\dd\Omega_I=0,\qquad \dd h_I=c_I\Omega_I,
\qquad c_A=-\delta q_Af_A,\quad c_+=-q_+f_+,\quad c_-=-q_-f_-.
\label{eq:hopf-curvatures}
\end{equation}
Thus $\Omega_A$ has norm squared $-\delta$, and the other area forms have norm squared $1$. For $f=2$, $\Omega_A=\theta^{01}$; for $f=0$, $\Omega_A=\theta^{12}$. Write precisely the following six flux components:
\begin{align}
F={}&\sum_{I=A,+,-}\mathcal E_I\,\Omega_I\wedge h_I\wedge\dd x
+\mathcal T_A\Omega_A\wedge h_+\wedge h_-\notag\\
&+\mathcal T_+h_A\wedge\Omega_+\wedge h_-
+\mathcal T_-h_A\wedge h_+\wedge\Omega_-.
\label{eq:six-flux}
\end{align}
The metric is a strict product of these Berger factors and the unwarped flat $x,y$ plane; all coefficients are constant. No other invariant four-form is admitted in this lemma.

Define $S_E=-\mathcal E_A^2+\mathcal E_+^2+\mathcal E_-^2$ and $S_T=-\mathcal T_A^2+\mathcal T_+^2+\mathcal T_-^2$. Since $F^2=24(S_E+\delta S_T)$, the two flat equations are
\begin{equation}
0=R_{yy}=-\frac{F^2}{144},\qquad
0=R_{xx}=\frac12S_E-\frac{F^2}{144}.
\label{eq:flat-equations}
\end{equation}
They imply $F^2=S_E=S_T=0$. For $\chi_A=-1$, $\chi_+=\chi_-=1$, the geometric Ricci eigenvalues and the flux stress eigenvalues are respectively
\begin{align}
\frac{R^{(I)}_{\rm base}}{g^{(I)}_{\rm base}}
&=\chi_If_I^2\left(1-\frac{q_I^2}{2}\right)
=\frac{\chi_I}{2}(\mathcal E_I^2+\delta\mathcal T_I^2),\notag\\
\frac{R^{(I)}_{\rm fibre}}{g^{(I)}_{\rm fibre}}
&=\frac{\chi_If_I^2q_I^2}{2}
=\frac{\chi_I}{2}(\mathcal E_I^2-\delta\mathcal T_I^2).
\label{eq:ricci}
\end{align}
The geometric expressions follow by inserting the rescaled structure constants $C^k{}_{ij}=C^{(0)k}{}_{ij}d_k/(d_id_j)$ in the Koszul connection. Adding and subtracting the two equations gives
\begin{equation}
\mathcal E_I^2=f_I^2,\qquad
\mathcal T_I^2=\delta f_I^2(1-q_I^2).
\label{eq:stress-solution}
\end{equation}
Closure of \eqref{eq:six-flux} has three independent coefficients:
\begin{equation}
\mathcal T_Ac_++\mathcal T_+c_A=0,\qquad
-\mathcal T_Ac_-+\mathcal T_-c_A=0,\qquad
-\mathcal T_+c_--\mathcal T_-c_+=0.
\label{eq:six-bianchi}
\end{equation}
Because every $c_I$ is nonzero, there is a real $\zeta$ such that
$\mathcal T_A/c_A=-\mathcal T_+/c_+=\mathcal T_-/c_-=\zeta$. Equations~\eqref{eq:stress-solution} then imply
\begin{equation}
q_A^2=q_+^2=q_-^2=\frac1{1+\delta\zeta^2},\qquad
f_A^2=f_+^2+f_-^2.
\label{eq:squash-necessity}
\end{equation}
The second relation follows from $S_E=0$. This proves necessity within \eqref{eq:six-flux}, including the undeformed branch. It does not yet prove the Maxwell equation, and it is not a statement about all invariant fluxes.

For existence, \eqref{eq:scaled} has $q_I=\lambda$, $\mathcal E_I=f_I$, and
\begin{equation}
\mathcal T_A=-\nu f_A,\qquad
\mathcal T_+=\delta\nu f_+,\qquad
\mathcal T_-=-\delta\nu f_-.
\label{eq:magnetic-amplitudes}
\end{equation}
Every pair of the six four-forms has a repeated index, hence $F\wedge F=0$. The electric Hodge-duals are closed. Differentiating the magnetic duals leaves a single coefficient,
\begin{equation}
\delta(-c_A\mathcal T_A-c_+\mathcal T_++c_-\mathcal T_-)
=-\lambda\nu(f_A^2-f_+^2-f_-^2)=0.
\label{eq:maxwell-coefficient}
\end{equation}
Together with \eqref{eq:ricci}, this verifies all equations of motion. The restriction to six forms is essential: Appendix~\ref{app:unequal} gives an unequal-sphere solution with nonzero $F^2$ and $F\wedge F$, including a point with exactly eight real supercharges.

\section{Exact supersymmetry and the exceptional algebra}\label{sec:susy}

The regular non-null solutions of Section~\ref{sec:geometry} preserve exactly eight local real supercharges whenever $\nu\ne0$; the undeformed solution preserves sixteen. We prove both statements directly from the eleven-dimensional Killing-spinor equation. We then determine the odd brackets, including their relative normalization, rather than identifying the superalgebra from the number of spinors alone. All statements in this section concern a simply connected local patch; descent to a quotient is a separate condition.

\subsection{The seed solution space}

We use the Clifford conventions of \eqref{eq:kse}, with the five-index gamma matrix fully antisymmetrized. For any four-form,
\begin{equation}\label{susy:slash}
 \mathscr D_i=\nabla_i+\frac1{12}\Gamma_i\slashed F
 -\frac14\slashed{\iota_iF},\qquad
 \slashed F=\frac1{4!}F_{jklm}\Gamma^{jklm},\qquad
 \slashed{\iota_iF}=\frac1{3!}F_{ijkl}\Gamma^{jkl}.
\end{equation}
In particular, $\slashed F_0=2HX$. On the seed, the flat-direction equations and the equations within a three-dimensional factor are
\begin{align}
 \mathscr D_x&=\partial_x+\frac13H,&
 \mathscr D_y&=\partial_y+\frac16YHX,\label{susy:seed-flat}\\
 \mathscr D_i&=\nabla_i+\Gamma_i\left(\frac16H-\frac r2V\right)X,
 & (r,V)&=(a,A),(b,B),(c,C),\label{susy:seed-factor}
\end{align}
where the last line applies to an index in the corresponding factor. The volume Clifford elements are parallel in the product metric. Using $\{H,V\}=2rV^2$ gives the mixed supercurvature
\begin{equation}\label{susy:seed-upper}
 [\mathscr D_y,\mathscr D_i]=\frac r{18}Y\Gamma_iVH.
\end{equation}
Since $r>0$ and $Y\Gamma_iV$ is invertible, every Killing spinor satisfies $H\epsilon=0$. No independence of $x$ or $y$ has been assumed: it follows from \eqref{susy:seed-flat} after imposing this necessary condition.

The radius relation \eqref{eq:radii} implies
\begin{equation}\label{susy:seed-projector}
 H^2=0,\qquad AH=a(1+Q),\qquad Q^2=1,\qquad
 \operatorname{tr}Q=0.
\end{equation}
Thus $\ker H=\operatorname{im}P_H$ has real dimension sixteen. On this subbundle the remaining equations become
\begin{equation}\label{susy:seed-reduced}
 \nabla_i\epsilon=\frac r2\Gamma_iVX\epsilon,
 \qquad \partial_x\epsilon=\partial_y\epsilon=0.
\end{equation}
They preserve $\ker H$. To see their integrability explicitly, put
$\chi_A=-1$ and $\chi_+=\chi_-=1$. Within factor $I$, the spin curvature is
$\chi_Ir_I^2\Gamma_{ij}/2$, whereas
\begin{equation}\label{susy:seed-flatness}
 \left[\frac{r_I}{2}\Gamma_iV_IX,
       \frac{r_I}{2}\Gamma_jV_IX\right]
 =-\frac{\chi_Ir_I^2}{2}\Gamma_{ij}.
\end{equation}
The mixed commutators between distinct factors vanish. Hence the restricted superconnection is flat, and each initial value in $\ker H$ extends uniquely to a local solution. This proves that the seed has exactly sixteen real Killing spinors. Since $X$ commutes with $Q$, its two eigenspaces split these into two eight-dimensional sectors. We shall transport
\begin{equation}\label{susy:seed-sector}
 W=\operatorname{im}\Pi_0,\qquad H\epsilon_0=0,\qquad X\epsilon_0=\epsilon_0.
\end{equation}
The real count uses the Majorana spinor bundle of signature $(1,10)$; the real bosonic superconnection and the projectors preserve its real structure.

\subsection{The full deformed connection and its kernel}

For the scale-removed fields in \eqref{eq:scaled}, write
$e^i=d_i\theta^i$, where $d_i=\lambda$ on the selected fibers $f,5,8$ and $d_i=1$ otherwise. Here $f=2$ for $\delta=1$ and $f=0$ for $\delta=-1$. Let $E_i$ be the dual vectors, and define structure constants by
$\dd e^k=-\tfrac12 C^k{}_{ij}e^i\wedge e^j$. Their actual values and the superconnection are
\begin{equation}\label{eq:connection}
\begin{aligned}
 C^k{}_{ij}&=C^{(0)k}{}_{ij}\frac{d_k}{d_id_j},
 &C^{(0)k}{}_{ij}&=2r_I\eta^{kk}\varepsilon_{kij}
 \quad (i,j,k\in I),\\
 \mathscr D_i&=E_i+\mathcal A_i,
 &\mathcal A_i&=\frac14\omega_{i ab}\Gamma^{ab}
       +\frac1{12}\Gamma_i\mathcal F
       -\frac1{24}\widetilde F_{iabc}\Gamma^{abc},\\
 \omega_{i ab}&=\frac12(C_{aib}-C_{iba}+C_{bai}),
 &\mathcal F=\slashed{\widetilde F}&=2H(X-\delta\nu P),
 \qquad P=\Gamma^{f58}.
\end{aligned}
\end{equation}
Here $C_{aib}=\eta_{ak}C^k{}_{ib}$, the structure constants vanish between distinct factors, and
$\omega_{i ab}=\langle E_a,\nabla_{E_i}E_b\rangle$ is the Koszul connection. Formula~\eqref{eq:connection}, together with the explicit four-form in \eqref{eq:scaled}, specifies all eleven differential operators. The overall constant rescaling $(g,F)=(s^2\widetilde g,s^3\widetilde F)$ multiplies their orthonormal components by $s^{-1}$ and does not alter their kernel.

The useful Clifford relations are
\begin{equation}\label{susy:deformed-projector}
 P^2=-\delta,\qquad [P,H]=0,\qquad \{P,X\}=0,
 \qquad J=(\lambda+\delta\nu P)X,
 \qquad \Pi=\frac14(1-Q)(1+J).
\end{equation}
Since $\lambda^2+\delta\nu^2=1$, one has $J^2=1$, $[J,Q]=0$ and $\operatorname{tr}\Pi=8$. These algebraic conditions become sufficient for the full differential equation through the invertible map
\begin{equation}\label{susy:spinor-map}
 S=\frac{1+\lambda+\delta\nu P}{\sqrt{2(1+\lambda)}},
 \qquad S^2=\lambda+\delta\nu P,\qquad
 SXS^{-1}=J,\qquad [S,H]=0.
\end{equation}
In fact, substitution in \eqref{eq:connection} gives, for every direction,
\begin{equation}\label{susy:transport}
 d_i\mathcal A_i S\Pi_0=S\mathcal A_i^{(0)}\Pi_0.
\end{equation}
Here $\mathcal A_i^{(0)}$ is obtained by setting $\lambda=1$, $\nu=0$ in the same connection. The identity follows using $H\Pi_0=0$, $X\Pi_0=\Pi_0$, $P^2=-\delta$, and $\lambda^2+\delta\nu^2=1$; it includes the flat directions. Importantly, the dual vectors also obey $E_i=d_i^{-1}E_i^{(0)}$. Since $S$ has constant components in these frames, \eqref{susy:transport} is the differential intertwining identity
\begin{equation}\label{susy:operator-transport}
 d_i\mathscr D_i(S\epsilon_0)=S\mathscr D_i^{(0)}\epsilon_0.
\end{equation}
Thus every seed solution in $W$ gives a deformed solution $\epsilon=S\epsilon_0$. This establishes eight independent local solutions without assuming that their frame components are constant.

To exclude additional solutions, the noncoordinate-frame curvature must include the structure-constant term:
\begin{equation}\label{susy:curvature-definition}
 \mathcal R_{ij}=[\mathcal A_i,\mathcal A_j]-C^k{}_{ij}\mathcal A_k.
\end{equation}
Two components of this operator suffice. Direct reduction of \eqref{eq:connection} yields
\begin{align}
 \mathcal R_{5y}&=-\frac{b\lambda^2}{18}Y\Gamma_5BH,
 \label{susy:upper-H}\\
 \mathcal R_{34}P_H&=
 \delta b^2\nu\,\Gamma^{34}P(\lambda-\delta\nu P)(1-J)P_H.
 \label{susy:upper-J}
\end{align}
The first identity forces $H\epsilon=0$ because $b\lambda>0$. In the second, $\Gamma^{34}$ and $P$ are invertible, and
\begin{equation}\label{susy:invertible}
 (\lambda-\delta\nu P)(\lambda+\delta\nu P)=1.
\end{equation}
For $\nu\ne0$, it therefore forces $J\epsilon=\epsilon$. The upper bound is precisely the eight-dimensional space already constructed. Consequently
\begin{equation}\label{susy:count}
 N_{\mathrm{local}}=8\quad(\nu\ne0),\qquad
 N_{\mathrm{local}}=16\quad(\nu=0).
\end{equation}
This argument applies to both signs of $\delta$ throughout the nondegenerate parameter domain. In particular, it excludes additional supersymmetry at special finite positive radius ratios.

\subsection{Normalized seed brackets}

We determine the seed bracket in the sector $W$ before transporting it. For a Majorana pairing $\bar\epsilon=\epsilon^T\mathcal C$, use
$(\Gamma^i)^T\mathcal C=-\mathcal C\Gamma^i$ and $\mathcal C^T=-\mathcal C$. If $M^\sharp=\mathcal C^{-1}M^T\mathcal C$, then
\begin{equation}\label{susy:pairing}
 \Pi_0^\sharp=\frac14(1+Q)(1-X),\qquad
 \Pi_0^\sharp A=A\Pi_0,\qquad
 \Omega(\epsilon,\chi)=\bar\epsilon A\chi.
\end{equation}
Since $A^2=1$, the last expression is a nondegenerate alternating pairing on $W$.

Let $K_j$ denote the surviving opposite-side invariant Killing fields, normalized by $K_j(o)=E_j^{(0)}(o)$ at a reference point. Their brackets and the radius-independent generators are
\begin{equation}\label{susy:even-normalization}
 [K_j,K_k]=-2r_I\eta^{\ell\ell}\varepsilon_{jk\ell}K_\ell,
 \qquad T_j=-\frac{K_j}{2r_I},\qquad j,k\in I.
\end{equation}
These $K_j$ are not globally the left-invariant frame vectors. At $o$,
$\nabla_a(K_j)_b=r_I\varepsilon_{jab}$. Combining this identity with \eqref{susy:seed-reduced} in the spinorial Lie derivative gives
\begin{equation}\label{susy:even-odd}
 \mathcal L_K\epsilon=K^i\nabla_i\epsilon+\frac14(\nabla_iK_j)\Gamma^{ij}\epsilon,
 \qquad
 \rho(T_j)\epsilon=-\frac12\Gamma_jV_I\epsilon.
\end{equation}
The three commuting sets of degree-two Clifford matrices preserve $W$ and $\Omega$, realizing the real form of $(\mathbf2,\mathbf2,\mathbf2)$. Their invariant alternating pairing is unique up to normalization.

The relation between this pairing and the physical vector current is
\begin{equation}\label{susy:compression}
 r_I\bar\epsilon A\Gamma_jV_I\chi
 =a\bar\epsilon\Gamma_j\chi,
 \qquad \epsilon,\chi\in W,
 \qquad (r_I,V_I)=(a,A),(b,B),(c,C).
\end{equation}
For the AdS factor this is immediate from $A^2=1$. A short projector argument proves the spherical identities. Put
$\mathsf S=AB$, $\mathsf T=AC$, $p=b/a$, $q=c/a$ and
$Q'= -p\mathsf S+q\mathsf T$. With
$P_-=(1-Q)/2$ and $P'_+=(1+Q')/2$, one has
\begin{equation}\label{susy:compression-proof}
 2pP'_+\mathsf SP_-=P'_+(Q-Q')P_-=-2P'_+P_-.
\end{equation}
For $j$ in the plus sphere, $Q\Gamma_j=\Gamma_jQ'$ and
$A\Gamma_jB=-\Gamma_j\mathsf S$. Inserting the projectors on both spinors in \eqref{susy:compression-proof} gives \eqref{susy:compression} with $r_I=b$. Interchanging $(\mathsf S,p)$ and $(\mathsf T,q)$ proves the minus-sphere identity. Both divisions are legitimate for finite positive radii.

Normalize the invariant form on each simple Lie algebra by a common trace convention,
\begin{equation}\label{susy:invariant-form}
 h_I(U,V)=-2\operatorname{tr}_{\mathbf2}\bigl(\rho(U)\rho(V)\bigr),
 \qquad h_I(T_j,T_k)=\chi_I\eta_{jk}.
\end{equation}
Define symmetric moment maps $\mu_I:W\times W\to\mathfrak g_I$ by
\begin{equation}\label{susy:moment-map}
 h_I\bigl(\mu_I(\epsilon,\chi),T_j\bigr)
 =\Omega\bigl(\epsilon,\rho(T_j)\chi\bigr).
\end{equation}
Writing $\kappa^j=\bar\epsilon\Gamma^j\chi$, equations
\eqref{susy:compression} and \eqref{susy:even-odd} imply
\begin{equation}\label{susy:moment-coefficients}
 \Omega\bigl(\epsilon,\rho(T_j)\chi\bigr)
 =-\frac{a}{2r_I}\kappa_j,
 \qquad
 \mu_I=-\frac{\chi_Ia}{2r_I}\kappa^jT_j.
\end{equation}
The flat components of the vector current vanish by the $Q$ and $X$ projections. Equivariance selects the three simple-factor Killing fields in \eqref{susy:even-normalization}; there is no invariant symmetric scalar pairing that could produce a spectator vector. Thus their normalized coefficients give
\begin{equation}\label{susy:seed-odd-bracket}
 K_0(\epsilon,\chi)
 =-\sum_I2r_I\kappa^jT_j
 =4a\left(-\mu_A+\frac{b^2}{a^2}\mu_+
                    +\frac{c^2}{a^2}\mu_-\right).
\end{equation}
The common factor is removed by choosing the invariant odd pairing
$\Omega_{\mathrm{std}}=4a\Omega$, whose moment maps are
$\mu_I^{\mathrm{std}}=4a\mu_I$. Relative to these normalized invariant tensors the odd-bracket coefficients are therefore
\begin{equation}\label{susy:exceptional-parameter}
 (-1,\gamma_+,\gamma_-),\qquad
 \gamma_+=\frac{b^2}{a^2}=\frac{L^2}{R_+^2},\qquad
 \gamma_-=\frac{c^2}{a^2}=\frac{L^2}{R_-^2},\qquad
 \alpha=\frac{\gamma_+}{\gamma_-}=\frac{R_-^2}{R_+^2}.
\end{equation}
Their sum vanishes by the radius relation, as required by the odd Jacobi identity, and none vanishes at finite positive radii. Together with \eqref{susy:even-odd}, these brackets define the real form of $\mathfrak d(2,1;\alpha)$ with even algebra
$\mathfrak{sl}(2,\mathbb R)\oplus\mathfrak{su}(2)_+\oplus\mathfrak{su}(2)_-$, in agreement with the seed superalgebra described in \cite{BachasEtAl2014}. Our convention defines $\alpha$ by the ratio of the two compact bracket coefficients; exchanging the spheres sends $\alpha$ to $\alpha^{-1}$. Thus the algebra parameter displayed in Eq.~(3.20) of Ref.~\cite{Noris2026} is $\alpha^{-1}$ in our convention.

\subsection{Transport of the odd brackets}

The spinor map \eqref{susy:spinor-map} also determines the deformed vector bilinears. Since $P^\sharp=P$, one has $S^\sharp=S$. The Clifford element $P$ anticommutes with a base gamma matrix and commutes with a fiber gamma matrix. Consequently
\begin{align}
 \overline{S\epsilon_1}\Gamma^iS\epsilon_2
 &=\bar\epsilon_1\Gamma^i\epsilon_2,
 &&i\text{ a base index},\label{susy:base-bilinear}\\
 \overline{S\epsilon_1}\Gamma^fS\epsilon_2
 &=\lambda\bar\epsilon_1\Gamma^f\epsilon_2
   +\delta\nu\bar\epsilon_1\Gamma^fP\epsilon_2,
 &&f\text{ a fiber index}.\label{susy:fiber-bilinear}
\end{align}
The additional term vanishes: $\Gamma^fP$ commutes with $X$, whereas $X\epsilon_j=\epsilon_j$ and $X^\sharp=-X$. Including the physical inverse frame, which scales by $s^{-1}$ on the base and $(s\lambda)^{-1}$ on the fibers, gives the equality of coordinate vector fields
\begin{equation}\label{susy:vector-transport}
 K'(S\epsilon_1,S\epsilon_2)=s^{-1}K_0(\epsilon_1,\epsilon_2).
\end{equation}
Thus the map $\epsilon\mapsto\sqrt{s}\,S\epsilon$ preserves the odd-odd bracket exactly. It is also equivariant under the surviving
$SL(2,\mathbb R)\times SU(2)_+\times SU(2)_-$ action: the coframe tensors defining $S$ are invariant under that action, and \eqref{susy:operator-transport} intertwines the corresponding spinor solutions. Hence the even-odd brackets are transported as well.

For every regular nonzero deformation, the complete odd space is the image of $W$. Its odd-generated Killing superalgebra is therefore a single $\mathfrak d(2,1;\alpha)$ with the unchanged parameter \eqref{susy:exceptional-parameter}. This statement distinguishes the algebra generated by the Killing spinors from an enlarged symmetry algebra that also includes bosonic spectator isometries. At $\nu=0$, the second seed spinor sector returns; the eight-supercharge conclusion is not asserted at that endpoint.

\section{A regular null family and its four supersymmetry branches}\label{sec:null}

The null deformation admits a nonsingular orthonormal frame and a direct Killing-spinor proof. With the electric flux orientation fixed, the two sphere Maurer--Cartan signs distinguish backgrounds with eight, four, or no local supercharges. We retain an arbitrary real amplitude $\beta$ throughout; no rescaling of a periodic coordinate is required.

\subsection{Geometry, flux, and field equations}

For $z>0$, define
\begin{align}
 e^0&=L\,\dd v+\frac L2(1+\beta^2)\frac{\dd u}{z^2},&
 e^1&=-L\,\dd v+\frac L2(1-\beta^2)\frac{\dd u}{z^2},&
 e^2&=L\frac{\dd z}{z},\label{null:frame}\\
 n&=e^0+e^1=Lz^{-2}\dd u,&
 \dd e^0&=a(1+\beta^2)n\wedge e^2,&
 \dd e^1&=a(1-\beta^2)n\wedge e^2.\nonumber
\end{align}
The two round spheres have independent signs $s_\pm\in\{1,-1\}$:
\begin{equation}\label{null:mc}
 \dd e^i=-s_+b\,\varepsilon_{ijk}e^{jk}\quad(i,j,k\in\{3,4,5\}),\qquad
 \dd e^i=-s_-c\,\varepsilon_{ijk}e^{jk}\quad(i,j,k\in\{6,7,8\}).
\end{equation}
Both repeated indices are summed, so $\dd e^5=-2s_+b e^{34}$ and $\dd e^8=-2s_-c e^{67}$. The complete fields are
\begin{align}
 g_N={}&L^2\left(\frac{\dd z^2-2\dd u\,\dd v}{z^2}-\frac{\beta^2\dd u^2}{z^4}\right)
       +g_++g_-+\dd x^2+\dd y^2,\label{null:metric}\\
 F_N={}&2a e^{0129}+2b e^{3459}+2c e^{6789}
 -2\beta n\wedge\left(ae^{258}+s_+b e^{348}-s_-c e^{567}\right)\nonumber\\
 ={}&F_0-\beta\,\dd(n\wedge e^5\wedge e^8).\label{null:flux}
\end{align}
Here $a=L^{-1}$, $b=R_+^{-1}$, $c=R_-^{-1}$, and $a^2=b^2+c^2$. In particular, the flux coefficients are $2a,2b,2c$. The signs in \eqref{null:mc} are relative to the positive electric coefficients in \eqref{null:flux}; changing a frame together with its flux components is not a change of this relative choice.

The $(-,-)$ branch matches the round-sphere boundary $\omega_{\rm N}=1$ of Eq.~(3.1) in Ref.~\cite{Noris2026}. Explicitly, set $u_{\rm N}=z$, $x_-=u$, $x_+=-2v$, $\alpha_{\rm null,N}=\beta^2$, $\beta=-\sqrt{\alpha_{\rm null,N}}$, and $A_0=n$, $A_\pm=2e^{5,8}$. With physical three-volumes aligned, $x=-z_1$ fixes the electric-leg ordering and $y=-z_2$ preserves the total orientation. The source's $\sqrt{\alpha_{\rm null,N}}\,\dd(A_0\wedge A_+\wedge A_-)/4$ is then the last term of \eqref{null:flux}. The $(+,+)$ branch changes the relative Hopf action; it is not an enhancement locus of this fixed-chirality source branch.

The field equations follow without a singular null-fiber limit. Bianchi is immediate from \eqref{null:flux}, and the wedge products give $F_N\wedge F_N=0$. The undeformed and additional fluxes are orthogonal, the latter is null, and hence $F_N^2=96(-a^2+b^2+c^2)=0$. More explicitly, exterior differentiation of the Hodge dual gives
\begin{equation}\label{null:maxwell}
\dd\star F_N=4\beta(-a^2+b^2+c^2)\,n\wedge e^{234679\,10}=0.
\end{equation}
For Einstein's equation, let $g_{N,A}$ be the deformed three-dimensional block in \eqref{null:metric}. Direct Cartan curvature gives
\begin{equation}\label{null:einstein}
 \operatorname{Ric}_{N,A}=-2a^2g_{N,A}+4a^2\beta^2 n\otimes n,
 \qquad \operatorname{Ric}_+=2b^2g_+,
 \qquad \operatorname{Ric}_-=2c^2g_-.
\end{equation}
The transverse three-form $ae^{258}+s_+be^{348}-s_-ce^{567}$ has norm squared $a^2+b^2+c^2=2a^2$. Its contribution to $F_{iPQR}F_j{}^{PQR}/12$ is therefore $4a^2\beta^2 n_i n_j$, precisely the extra term in \eqref{null:einstein}; the mixed electric--null contractions vanish. Thus all four sign choices solve the same bosonic equations.

\subsection{A finite set of supercurvature identities}

We use $\mathscr D_i=E_i+\mathcal A_i$ and the general Koszul and flux terms of \eqref{eq:connection}, now with the frame \eqref{null:frame}--\eqref{null:mc}. Its special non-null expression for $\mathcal F$ is replaced by \eqref{null:slash} below. Introduce
\begin{equation}\label{null:operators}
\begin{gathered}
 U=\Gamma^0+\Gamma^1,\qquad V=\Gamma^0-\Gamma^1,\qquad
 \Lambda=\Gamma^{01},\\
 P_U=\frac{1+\Lambda}{2},\qquad \mathcal N=U\Gamma^{58},\qquad
 \Delta=b^2(1-s_+)+c^2(1-s_-).
\end{gathered}
\end{equation}
Then $U^2=\mathcal N^2=0$, $\ker U=\operatorname{im}P_U$, $[H,\Lambda]=0$, $[Q,\mathcal N]=0$, and $\{\mathcal N,X\}=0$. In particular $VU=-2(1-\Lambda)$. Explicitly, the spin-connection matrices are
\begin{align}
 \Omega_0&=\frac a2(\Gamma^{02}+\beta^2U\Gamma^2),&
 \Omega_1&=\frac a2(-\Gamma^{12}+\beta^2U\Gamma^2),&
 \Omega_2&=-\frac a2\Lambda,\nonumber\\
 \Omega_i&=-\frac{s_+b}{2}\Gamma^iB\ (i=3,4,5),&
 \Omega_i&=-\frac{s_-c}{2}\Gamma^iC\ (i=6,7,8),&
 \Omega_x&=\Omega_y=0.\label{null:spinconnection}
\end{align}
A useful short form of the flux is
\begin{equation}\label{null:slash}
 \slashed F_N=2HX-2\beta U\mathcal W,\qquad
 \mathcal W=a\Gamma^{258}+s_+b\Gamma^{348}-s_-c\Gamma^{567}.
\end{equation}
For example, if $\mathcal W_i$ denotes contraction of this transverse three-form, then
$\slashed{\iota_iF_d}=-2\beta[(\delta_{i0}+\delta_{i1})\mathcal W-U\mathcal W_i]$.
These expressions determine every $\mathcal A_i$ without a choice of explicit gamma matrices. The curvature is the non-coordinate one,
$\mathcal R_{ij}=[\mathcal A_i,\mathcal A_j]-C^k{}_{ij}\mathcal A_k$.

Expanding these short expressions and using $a^2=b^2+c^2$ gives
\begin{align}
 U\mathcal R_{5y}&=-\frac b{18}UY\Gamma^5BH,&
 \mathcal R_{xy}P_U&=0,\nonumber\\
 \mathcal R_{xy}P_H&=\frac{\beta\Delta}{6}U\Gamma^{58}YP_H,&
 \mathcal R_{5y}P_U&=-\frac b{18}Y\Gamma^5BH P_U.\label{null:universal-curvature}
\end{align}
These are operator identities, not equations obtained by assuming the desired projectors. All Clifford reductions use the fixed irreducible convention $\Gamma^{0\cdots10}=-1$.

Suppose first that $\Delta>0$, which includes exactly the three sign choices other than $(+,+)$. The first identity in \eqref{null:universal-curvature} implies $UH\epsilon=0$, since $Y\Gamma^5B=Y\Gamma^{34}$ is invertible and anticommutes with $U$. To avoid an unjustified replacement of $UH$ by $H$, decompose $\epsilon=\epsilon_++\epsilon_-$ with $\Lambda\epsilon_\pm=\pm\epsilon_\pm$. The map $U$ is injective on the negative eigenspace, and $H$ preserves this decomposition. Consequently $H\epsilon_-=0$. The middle two identities then give
\begin{equation}\label{null:negative-sector}
 0=\mathcal R_{xy}\epsilon=\mathcal R_{xy}P_H\epsilon_-
   =\frac{\beta\Delta}{6}U\Gamma^{58}Y\epsilon_-.
\end{equation}
For $\beta\ne0$, invertibility of $\Gamma^{58}Y$ within that eigenspace forces $\epsilon_-=0$. Thus $U\epsilon=0$, and only now does the last identity in \eqref{null:universal-curvature} imply $H\epsilon=0$. Every solution in these three branches lies in the eight-dimensional space $\operatorname{im}P$, where $P=P_HP_U$.

\subsection{The four-supercharge branch and the mixed-sign obstruction}

For $(s_+,s_-)=(-,-)$, one further factorization is decisive:
\begin{equation}\label{null:four-upper}
 \mathcal R_{0x}P=\frac{a^2\beta}{3}\Gamma^{58}(1+X)P.
\end{equation}
Its prefactor is invertible for nonzero $\beta$, so every solution satisfies
\begin{equation}\label{null:four-projector}
 H\epsilon=0,\qquad U\epsilon=0,\qquad X\epsilon=-\epsilon,
 \qquad \Pi_4=\frac18(1-Q)(1+\Lambda)(1-X).
\end{equation}
The three commuting projector factors have trace $4$ in combination. This gives an upper bound, not yet existence. The latter follows from the explicit restricted connection:
\begin{equation}\label{null:four-connection}
 \begin{aligned}
 \mathcal A_{0,1,2,x,y}\Pi_4&=0,\\
 \mathcal A_i\Pi_4&=b\Gamma^iB\Pi_4\quad(i=3,4,5),&
 \mathcal A_i\Pi_4&=c\Gamma^iC\Pi_4\quad(i=6,7,8).
 \end{aligned}
\end{equation}
For clarity, the terms linear in $\beta$ in the first line vanish because
$(\mathcal W-\Gamma^{58}H)P_U=0$ in this branch. All quadratic terms contain $U\Gamma^2$. In sphere directions, $U\mathcal W_i$ vanishes on the projector; the spin and electric-flux terms each contribute half of the displayed answer.

The remaining matrices preserve all three constraints. Their flatness is elementary: on the first sphere,
$\mathcal A_3=b\Gamma^{45}$, $\mathcal A_4=-b\Gamma^{35}$, and $\mathcal A_5=b\Gamma^{34}$, so
\begin{equation}\label{null:four-flatness}
 [\mathcal A_3,\mathcal A_4]=-2b^2\Gamma^{34}
 =C^5{}_{34}\mathcal A_5,\qquad C^5{}_{34}=-2b.
\end{equation}
Cyclic permutations and the second sphere are identical, and different factors commute. Thus $\Pi_4$ defines a parallel flat subbundle. Every allowed initial value extends uniquely on a simply connected patch, proving exactly four solutions.

For mixed signs, the same preliminary elimination gives $\epsilon=P\epsilon$, but now
\begin{equation}\label{null:mixed-obstruction}
 H\Gamma^{58}Y\mathcal R_{0y}P
 =-\frac{(s_+-s_-)abc\beta}{3}XY P.
\end{equation}
Left multiplication of a curvature equation is legitimate without assuming the multiplier invertible. The right-hand side is invertible on the candidate space: $s_+-s_-=\pm2$, $abc\beta\ne0$, and $(XY)^2=-1$. Hence $\epsilon=0$ in both mixed branches. No exceptional positive radius ratio escapes this argument.

\subsection{The eight-supercharge branch}

For $(s_+,s_-)=(+,+)$, the correct first equation is stronger than its $U$-multiple:
\begin{equation}\label{null:eight-upper}
 \mathcal R_{5y}=-\frac b{18}Y\Gamma^5BH,\qquad
 \mathcal R_{02}P_H=\frac{a^2\beta}{2}\Gamma^{258}(1+2\Lambda)
 (1+\beta\mathcal N)(1-J_N)P_H,
 \quad J_N=(1-\beta\mathcal N)X.
\end{equation}
The first equality forces $H\epsilon=0$, but not $U\epsilon=0$. In the second equality all prefactors are invertible:
$(\Gamma^{258})^2=-1$, $(1+2\Lambda)^{-1}=(2\Lambda-1)/3$, and
$(1+\beta\mathcal N)^{-1}=1-\beta\mathcal N$. Hence $J_N\epsilon=\epsilon$ for $\beta\ne0$. Since $J_N^2=1$ and $[Q,J_N]=0$, the candidate projector
\begin{equation}\label{null:eight-projector}
 \Pi_8=\frac14(1-Q)(1+J_N)
 =S_N\Pi_0S_N^{-1},\qquad
 S_N=1-\frac\beta2\mathcal N
\end{equation}
has rank eight.

Existence is again visible in a small connection table. Direct substitution of \eqref{null:spinconnection}--\eqref{null:slash} gives
\begin{equation}\label{null:eight-connection}
 \mathcal A_iS_N\Pi_0=S_N M_i\Pi_0,\qquad
 \begin{aligned}
 M_0=M_1&=\frac a2V\Gamma^2,& M_2&=-a\Lambda,\\
 M_i&=-b\Gamma^iB\ (i=3,4,5),&
 M_i&=-c\Gamma^iC\ (i=6,7,8),\\
 M_x=M_y&=0.
 \end{aligned}
\end{equation}
Each $M_i$ preserves $\Pi_0$. The sphere connections are flat by the calculation in \eqref{null:four-flatness}, with both Maurer--Cartan and connection signs reversed. In the AdS block,
\begin{equation}\label{null:eight-flatness}
 [M_0,M_2]=-2aM_0
 =-a(1+\beta^2)M_0-a(1-\beta^2)M_1
 =C^k{}_{02}M_k.
\end{equation}
The $(1,2)$ equation is the same and the $(0,1)$ equation vanishes. Thus the restricted connection is flat in the deformed frame itself: the equality $M_0=M_1$ removes the apparent $\beta^2$ obstruction. This supplies eight solutions and completes the upper-bound argument.

The real structure in Appendix~\ref{app:clifford} commutes with these connections and projectors, so their complex ranks equal the dimensions of their Majorana real forms. The complete local count is therefore
\begin{equation}\label{null:count}
\begin{array}{c|cccc}
 (s_+,s_-)&(+,+)&(-,-)&(+,-)&(-,+)\\\hline
 N_{\rm local}\ (\beta\ne0)&8&4&0&0.
\end{array}
\end{equation}
At $\beta=0$ the curvature factors responsible for the additional projections vanish, and every branch returns to the sixteen-supercharge seed. The proof assumes $a,b,c>0$; decompactifying sphere limits are outside this assertion.

\subsection{Odd brackets and spectator symmetries}

The map $S_N$ in \eqref{null:eight-projector} transports the positive-$X$ seed Killing spinors. It also preserves their physical vector brackets. To make the frame dependence explicit, let $K_0^i=\bar\epsilon_{01}\Gamma^i\epsilon_{02}$ and let $K_N^i$ use $S_N\epsilon_{0r}$. The linear terms in $\beta$ vanish by the seed $X$ projection, while the quadratic terms give
\begin{equation}\label{null:vector-transport}
 \begin{pmatrix}K_N^0\\K_N^1\end{pmatrix}
 =\begin{pmatrix}1+\beta^2/2&\beta^2/2\\-\beta^2/2&1-\beta^2/2\end{pmatrix}
 \begin{pmatrix}K_0^0\\K_0^1\end{pmatrix},\qquad
 K_N^i=K_0^i\quad(i\ge2).
\end{equation}
The matrix displayed here is exactly the change of coframe from $\beta=0$ to \eqref{null:frame}; its inverse occurs in the dual vectors. Consequently $K_N^iE_i^N=K_0^iE_i^0$. The odd-generated algebra in the $(+,+)$ branch is the single real $\mathfrak d(2,1;\alpha)$ factor with the unchanged parameter \eqref{susy:exceptional-parameter}.

The four-supercharge branch has a different odd bracket. In the explicit Majorana realization of Appendix~\ref{app:clifford}, choose $\mathcal C=B_M\Gamma^0$. Since $B_M$ is real and symmetric, $\bar\epsilon_1\Gamma^i\epsilon_2=\epsilon_1^\dagger\Gamma^0\Gamma^i\epsilon_2$. The Hermitian involution $\Lambda$ is $+1$ on $\Pi_4$; for $i\ge2$, $\Gamma^0\Gamma^i$ anticommutes with $\Lambda$. Thus
\begin{equation}\label{null:four-bilinear}
 K^0=-\epsilon_1^\dagger\epsilon_2,
 \qquad K^1=\epsilon_1^\dagger\epsilon_2,
 \qquad K^i=0\quad(i\ge2).
\end{equation}
Since $E_0-E_1=L^{-1}\partial_v$, every odd bracket is along the central null translation $\mathsf N=\partial_v$. The pairing is nondegenerate: on Majorana spinors $\epsilon^\dagger\epsilon$ is a positive real norm. This establishes nondegeneracy without a numerical bilinear rank test.

The two full sphere $SU(2)$ actions have generators proportional to $\Gamma^iB/2$ and $\Gamma^mC/2$ on this space; each quadratic Casimir is $-3/4$. The four real spinors consequently carry the real form of $(\mathbf2,\mathbf2)$. With a common odd normalization their bracket can be written
\begin{equation}\label{null:four-algebra}
 \{Q_{a\dot a},Q_{b\dot b}\}
   =\varepsilon_{ab}\varepsilon_{\dot a\dot b}\mathsf N,
 \qquad[\mathsf N,Q]=0,
 \qquad
 (\mathfrak{su}(2)_+\oplus\mathfrak{su}(2)_-)\ltimes\mathfrak h_{1|4}.
\end{equation}
Here the final expression includes the sphere automorphisms; the odd generators alone generate only $\mathsf N$. Explicitly, take $\varepsilon_{12}=\varepsilon_{\dot1\dot2}=1$ and $(Q_{a\dot a})^*=-\varepsilon_{ab}\varepsilon_{\dot a\dot b}Q_{b\dot b}$. On this real form the pairing is negative definite: for coefficients satisfying the same reality condition, $\varepsilon_{ab}\varepsilon_{\dot a\dot b}z^{a\dot a}z^{b\dot b}=-2(|z^{1\dot1}|^2+|z^{1\dot2}|^2)$. Thus the sign agrees with \eqref{null:four-bilinear}; only a positive overall factor is absorbed in the real odd normalization. In particular this algebra is not $\mathfrak d(2,1;\alpha)$, whose odd part has dimension eight.

The remaining null $\mathfrak{sl}(2,\mathbb R)$, two Hopf symmetries, and flat translations are even spectators on these four spinors. One can establish their trivial action without integrating a spinorial Lie derivative: they preserve the fields and commute with the two full sphere $SU(2)$ actions and with $\mathsf N$. The real $(\mathbf2,\mathbf2)$ representation is absolutely irreducible, so Schur's lemma makes any such action a real scalar $r$. The Jacobi identity with the nondegenerate bracket \eqref{null:four-algebra} then gives $2r=0$. These are local statements; a quotient and its spin structure may reduce the number of globally defined supercharges.

% !TeX root = paper.tex
\section{Flux charges, compact quotients and global restrictions}\label{sec:charges}
\subsection{Which charges are defined?}
Fix the normalization
\begin{equation}
\begin{gathered}
2\kappa_{11}^2=(2\pi)^8\ell_p^9=16\pi G_{11},\\
N_5=\frac{1}{(2\pi\ell_p)^3}\int_{\Sigma_4}F,\qquad
Q_2^{\rm Page}=\frac{1}{(2\pi\ell_p)^6}\int_{\Sigma_7}
\left(\star F+\frac12C_3\wedge F\right).
\end{gathered}
\label{charge:definitions}
\end{equation}
The distinction between Maxwell, Page and source charge matters in a theory with a Chern--Simons term~\cite{Marolf2000}. Locally, $F=\dd C_3$ and \eqref{eq:eom} give $\dd(\star F+C_3\wedge F/2)=0$. Page charge can nevertheless depend on a large-gauge representative. Maxwell flux is gauge invariant but is not generally closed when $F\wedge F\ne0$.

On the internal space stated in the seed,
\begin{equation}
H_4(S^3_+\times S^3_-\times\mathbb R^2;\mathbb Z)
=H_7(S^3_+\times S^3_-\times\mathbb R^2;\mathbb Z)=0.
\end{equation}
Thus the compact charge cycles used below are absent before an additional quotient. For noncompact $x$, the seed has the global potential $C_3=-x(f_Av_A+f_+v_++f_-v_-)$. Integrating over $S^3_\pm\times I_x$, or the corresponding electric open chain, defines regulated flux densities, not integer brane charges. In the orientation specified below they are
\begin{equation}
\rho_{5,\pm}=\frac{R_\pm^2}{2\pi\ell_p^3},\qquad
\rho_2=\frac{R_+^3R_-^3}{8\pi^2L\ell_p^6}.
\end{equation}
The electric expression is a Maxwell density; an open-chain Page integral also needs endpoint gauge data.

\subsection{Torus periods, orientations and deformation invariance}
Now impose, as an additional choice,
\begin{equation}
x\sim x+\ell_x,\qquad y\sim y+\ell_y.
\label{charge:quotient}
\end{equation}
These coordinate periods are the proper seed circumferences. Their proper values in \eqref{eq:non-null} are $s\ell_x,s\ell_y$. The fields descend because they are independent of $x,y$. The Killing spinors also descend with the periodic spin structure: the seed KSE forces $\partial_x\epsilon_0=\partial_y\epsilon_0=0$, the non-null map is independent of $x,y$, and both reduced null connections have vanishing flat components. Therefore this particular torus quotient preserves the stated local supersymmetries; arbitrary further quotients need not do so.

Orient $\Sigma_{4,\pm}=S^3_\pm\times S^1_x$ by $v_\pm\wedge\dd x$. Since
\begin{equation}
\star e^{0129}=-e^{345678\,10},
\end{equation}
choose $\Sigma_{7,y}=S^3_+\times S^3_-\times S^1_y$ with orientation $-v_+\wedge v_-\wedge\dd y$. The seed integrals are then
\begin{align}
\int_{\Sigma_{4,\pm}}F_0&=\frac{2}{R_\pm}(2\pi^2R_\pm^3)\ell_x,\notag\\
\int_{\Sigma_{7,y}}\star F_0&=\frac{2}{L}(2\pi^2R_+^3)(2\pi^2R_-^3)\ell_y.
\end{align}
For the displayed deformations, the inclusion $\iota:\Sigma_{7,y}\hookrightarrow M$ obeys
\begin{equation}
\iota^*F=0,\qquad
\iota^*(C_3\wedge F)=0.
\label{charge:page-reduction}
\end{equation}
The second equality holds in every potential patch; the same pullback kills a large-gauge shift on this specified cycle. Consequently its Page integral equals its Maxwell integral. This argument does not make arbitrary Page charges gauge independent.

The resulting classical cycle integrals, in the natural normalization \eqref{eq:non-null}, are
\begin{equation}
\boxed{N_{5,\pm}=\frac{R_\pm^2\ell_x}{2\pi\ell_p^3},\qquad
N_2=\frac{R_+^3R_-^3\ell_y}{8\pi^2L\ell_p^6}.}
\label{charge:values}
\end{equation}
To establish their invariance, it is insufficient to observe that the added four-form is exact. Its magnetic pullback vanishes, but the electric integral also involves the changed Hodge star. In the seed coframe the relevant weights are
\begin{equation}
\frac{\operatorname{vol}(g')}{\operatorname{vol}(g_0)}=G^{-1/3},\qquad
(g'^{-1})^{\otimes4}\big|_{v_A\wedge\dd x}=G^{1/3}.
\label{charge:hodge-weights}
\end{equation}
The original electric term has one fibre and three base indices, so the weights cancel. Every term in $\star'\dd\Theta$ retains an AdS leg and vanishes on $\Sigma_{7,y}$. This proves \eqref{charge:values} at fixed seed radii and coordinate periods. For the null family, the $(u,v)$ metric block has determinant $-L^4/z^4$, independent of the amplitude. The original AdS volume and its internal dual are unchanged, while the added dual flux again retains an AdS leg. The same charges therefore apply to all its chirality branches.

By contrast, an independent overall scaling at fixed integration domain gives
\begin{equation}
(g,F,C_3)\mapsto(s^2g,s^3F,s^3C_3),\qquad
N_5\mapsto s^3N_5,\qquad N_2\mapsto s^6N_2.
\label{charge:scaling}
\end{equation}
Here $\star_{s^2g}|_{\Omega^4}=s^3\star_g$. One must not assign the invariant charges of \eqref{eq:non-null} directly to its scale-removed fields \eqref{eq:scaled} with the same proper periods.

\subsection{Quantization and the seed central charge}
Quantization is an additional condition on \eqref{charge:values}, not an automatic consequence of compactification. Impose $N_{5,+},N_{5,-},N_2\in\mathbb Z_{>0}$. The tangent bundles of the spheres and torus are trivial, so the gravitational shift $p_1/4$ in the magnetic flux-quantization condition~\cite{Witten1996flux} vanishes on the cycles in question. Nonzero magnetic charge still requires patchwise potentials.

The two magnetic periods share the same $\ell_x$. Hence
\begin{equation}
\frac{R_+^2}{R_-^2}=\frac{N_{5,+}}{N_{5,-}}\in\mathbb Q_{>0},\qquad
\alpha=\frac{N_{5,-}}{N_{5,+}}.
\label{charge:rationality}
\end{equation}
If the squared-radius ratio is irrational, no nonzero common $\ell_x$ makes both magnetic charges integral. Conversely, for $R_+^2/R_-^2=m/n$, one may choose $\ell_x=2\pi\ell_p^3m/R_+^2=2\pi\ell_p^3n/R_-^2$, then choose $\ell_y$ to obtain an integer electric charge. The continuous deformations above preserve these selected flux integers.

As a check confined to the undeformed AdS background, the internal volume and effective Newton constant are
\begin{equation}
V_8=4\pi^4R_+^3R_-^3\ell_x\ell_y,\qquad
G_{11}=16\pi^7\ell_p^9,\qquad G_3=G_{11}/V_8.
\end{equation}
With ordinary AdS Brown--Henneaux boundary conditions~\cite{BrownHenneaux1986},
\begin{equation}
c_L=c_R=\frac{3L}{2G_3}
=\frac{6N_2N_{5,+}N_{5,-}}{N_{5,+}+N_{5,-}}.
\label{charge:seed-central}
\end{equation}
This is a seed central charge. It is not a derivation of a warped asymptotic algebra. Without \eqref{charge:quotient}, the internal volume is infinite and this finite central charge is replaced, at most, by a regulated density.

\subsection{Causality and the boundary of the local analysis}
Local supersymmetry does not ensure a causal global completion. For the timelike branch, use a connection regular at the origin of the hyperbolic base,
\begin{equation}
\dd s_A^2=\frac{s^2L^2}{4}\left[
\dd\rho^2+\sinh^2\rho\,\dd\phi^2
-\lambda^2\big(\dd\tau+(\cosh\rho-1)\dd\phi\big)^2\right],
\qquad\phi\sim\phi+2\pi.
\end{equation}
At fixed $\tau$, the closed angular circles have
\begin{equation}
g_{\phi\phi}=\frac{s^2L^2}{4}(\cosh\rho-1)
\big[(1-\lambda^2)\cosh\rho+1+\lambda^2\big].
\end{equation}
For every nonzero timelike deformation, $\lambda>1$, this is negative when
$\cosh\rho>(\lambda^2+1)/(\lambda^2-1)$. Unwrapping the time coordinate does not remove these base-angle circles. Different global constructions require their own causal and spinorial analysis.

Finally, the local fields do not specify warped energy or angular momentum. Those observables require global identifications, integration surfaces, boundary falloffs, normalized generators, a boundary action and a reference zero. Brane-flux charges and the local Killing superalgebra do not supply the missing data. Thus the unspecified gravitational charges are not assigned the value zero, nor is \eqref{charge:seed-central} exported to the warped geometries.


\section{Conclusions}
The three-fibre sector admits a uniform description of spacelike and timelike deformations. The metric anisotropy is supported by an exact addition to the four-form, with the radius relation inherited from the undeformed product. Within the six-component flux ansatz, the Einstein equations and the Bianchi identity force all three fibre-to-base ratios to coincide. The unequal-sphere configuration of Appendix~\ref{app:unequal}, including an exact eight-supercharge point, demonstrates the boundary of that statement.

For every regular nonzero non-null deformation, two supercurvature factorizations impose precisely the two commuting projections needed for eight real Killing spinors. A connection-level spinor map supplies the matching lower bound. The normalized seed moment maps and the transported vector bilinear then determine a single $\mathfrak d(2,1;\alpha)$ factor with $\alpha=R_-^2/R_+^2$, rather than identifying the algebra by dimension alone. In the null sector, the relative Hopf chiralities are substantive data: the four branches have respectively eight, four, zero and zero real supercharges. The four-supercharge branch has a null supertranslation bracket, not the exceptional algebra.

The specified brane-flux charges are unchanged along the naturally normalized deformations at fixed seed periods. Compactification additionally requires quantization and a periodic spin structure; the resulting rationality constraint on the squared-radius ratio is not implied by the local equations. Local supersymmetry likewise does not remove the closed timelike curves of the standard timelike completion. These results separate local construction from global admissibility.

The reconstruction therefore stops at a definite boundary. It does not classify all invariant four-forms or all discrete quotients, and it does not replace the broader families of Ref.~\cite{Noris2026}. Determining asymptotic charges requires a chosen global spacetime, boundary conditions, normalized symmetry generators and a reference prescription. Those data constitute a separate problem, not an additional number fixed by the local Killing superalgebra.

\appendix
% !TeX root = paper.tex
\section{An unequal-sphere supersymmetric example}\label{app:unequal}
The necessity argument in Section~\ref{sec:geometry} excludes magnetic components that can change both the flat Einstein equations and the Chern--Simons contribution. The following explicit example exhibits this limitation. It is not a counterexample to a theorem whose hypotheses already restrict the flux to \eqref{eq:six-flux}.

For this appendix, reverse the Maurer--Cartan convention:
\begin{equation}
C^{(0)i}{}_{jk}=-f_I\eta^{ii}\varepsilon_{ijk}.
\label{unequal:mc}
\end{equation}
The selected spacelike fibres then obey $\dd\theta^2=f_A\theta^{01}$, $\dd\theta^5=f_+\theta^{34}$ and $\dd\theta^8=f_-\theta^{67}$. For $q^2+\nu^2=1$, $q>0$, scale only two fibres,
\begin{equation}
e^2=q\theta^2,\qquad e^5=q\theta^5,\qquad e^8=\theta^8,
\qquad q_A=q_+=q,\quad q_-=1,
\end{equation}
with all other coframe components unchanged. The four-form is
\begin{align}
F_{\rm u}={}&f_Ae^{0129}+f_+e^{3459}+qf_-e^{6789}\notag\\
&+\nu f_Ae^{015\,10}-\nu f_+e^{234\,10}+\nu f_-e^{0134}.
\label{unequal:flux}
\end{align}
In particular, the last term lies outside \eqref{eq:six-flux}.

\subsection{Field equations}
Set $E=qf_-$, $T_A=\nu f_A$, $T_+=-\nu f_+$, $M=\nu f_-$, $c_A=qf_A$ and $c_+=qf_+$. The independent Bianchi and Maxwell coefficients reduce to
\begin{align}
\dd F=0:&\qquad c_+T_A+c_AT_+=0,\notag\\
\dd\star F+\tfrac12F\wedge F=0:&\qquad
EM-c_AT_A-c_+T_+=0,\notag\\
&\qquad -c_AM+ET_A=0,\qquad c_+M+ET_+=0.
\label{unequal:form-equations}
\end{align}
The first Maxwell coefficient is $q\nu(f_-^2-f_A^2+f_+^2)$, and the others vanish identically. Here $F^2=-72\nu^2f_-^2$ and $F\wedge F\ne0$ for $\nu\ne0$. The Chern--Simons term supplies the three terms $EM$, $ET_A$ and $ET_+$ in \eqref{unequal:form-equations}; it cannot be discarded.

Retaining the trace term in \eqref{eq:eom}, the stress tensor gives
\begin{equation}
\begin{array}{c|cc}
 &R_{\rm base}/g_{\rm base}&R_{\rm fibre}/g_{\rm fibre}\\ \hline
\AdS_3&-f_A^2(1-q^2/2)&-f_A^2q^2/2\\
S^3_+&f_+^2(1-q^2/2)&f_+^2q^2/2\\
S^3_-&f_-^2/2&f_-^2/2
\end{array}
\end{equation}
and $R_{xx}=R_{yy}=0$. Distinct flux components share at most two indices, so the off-diagonal stress vanishes. These are exactly the curvatures of the stated product metric, proving the full bosonic equations for this family.

\subsection{A finite exact supersymmetry certificate}
A single supersymmetric point is sufficient to rule out an unrestricted equality claim. Choose a length unit and take
\begin{equation}
(f_A,f_+,f_-)=(10,6,8),\qquad (q,\nu)=(3/5,4/5).
\label{unequal:point}
\end{equation}
Define
\begin{gather}
P_{\rm u}=\Gamma^{2,5,10},\qquad
H_{\rm u}=5A+3B+4(q+\nu P_{\rm u})C,\notag\\
Q_{\rm u}=\frac35AB+\frac45A(q+\nu P_{\rm u})C,\qquad
J_{\rm u}=(q+\nu P_{\rm u})X,\qquad
\Pi_{\rm u}=\frac14(1-Q_{\rm u})(1-J_{\rm u}).
\label{unequal:projector}
\end{gather}
Clifford multiplication gives $Q_{\rm u}^2=J_{\rm u}^2=1$, $[Q_{\rm u},J_{\rm u}]=0$, and $\operatorname{tr}\Pi_{\rm u}=8$. The projected conditions are $H_{\rm u}\epsilon=0$ and $J_{\rm u}\epsilon=-\epsilon$.

For completeness, the finite matrix certificate can be specified without an unidentified rank computation. Construct the eleven connection matrices by \eqref{eq:kse} and the Koszul formula in \eqref{eq:connection}, using \eqref{unequal:mc}, $d_2=d_5=3/5$, all other $d_i=1$, and the following ordered four-form coefficients:
\begin{equation}
(F_{0129},F_{3459},F_{6789},F_{015\,10},F_{234\,10},F_{0134})
=(10,6,24/5,8,-24/5,32/5).
\label{unequal:matrix-input}
\end{equation}
Do not use the special equal-fibre expression for $\mathcal F$ in \eqref{eq:connection}; here $\mathcal F=\slashed F_{\rm u}$ is determined by \eqref{unequal:matrix-input}. With $\mathscr D_i=E_i+\mathcal A_i$, the matrices obey
\begin{equation}
(1-\Pi_{\rm u})\mathcal A_i\Pi_{\rm u}=0\quad(0\leq i\leq10),
\qquad \mathcal R_{ij}\Pi_{\rm u}=0\quad(i<j).
\label{unequal:flatness}
\end{equation}
These are finite rational Clifford identities with the explicitly specified operands. They show that the rank-eight subbundle is preserved and flat, giving eight local solutions, rather than merely eight vectors in a curvature kernel.

An explicit nonzero minor supplies the converse. In the tensor-product representation of Appendix~\ref{app:clifford}, order matrix rows and columns from $1$ to $32$ and set
\begin{equation}
I=\{1,\ldots,12,17,\ldots,28\},\qquad
\det\big[(\mathcal R_{01})_{I,I}\big]\equiv627\pmod{1009}.
\label{unequal:minor}
\end{equation}
For this modular evaluation all rational denominators are invertible and $i$ may be represented by $469$, since $469^2=-1\pmod{1009}$. The nonzero $24\times24$ minor proves $\operatorname{rank}\mathcal R_{01}\geq24$, so no more than eight initial values are possible. Together with \eqref{unequal:flatness} and the Majorana structure below, the point \eqref{unequal:point} has exactly eight real supercharges. The first sphere has fibre-to-base ratio $3/5$ while the second is round. Their distinct Ricci eigenvalue ratios make the inequality geometric, not a coframe artifact. This certificate does not assert a supersymmetry classification of the entire unequal-sphere family.

\section{Clifford representation and real spinors}\label{app:clifford}
A fixed matrix realization makes the finite certificate above reproducible and clarifies real versus complex dimensions. Let $\sigma_1,\sigma_2,\sigma_3$ be the usual Pauli matrices. Define ten Euclidean matrices by
\begin{align}
\widetilde\Gamma^{2j}&=\sigma_3^{\otimes j}\otimes\sigma_1
\otimes\mathbf1_2^{\otimes(4-j)},\notag\\
\widetilde\Gamma^{2j+1}&=\sigma_3^{\otimes j}\otimes\sigma_2
\otimes\mathbf1_2^{\otimes(4-j)},\qquad j=0,\ldots,4.
\end{align}
Take $\Gamma^0=i\widetilde\Gamma^0$, $\Gamma^i=\widetilde\Gamma^i$ for $1\leq i\leq9$, and $\Gamma^{10}=\sigma_3^{\otimes5}$. Tensor factors and their two-dimensional basis vectors are ordered as written. Then
\begin{equation}
\{\Gamma^i,\Gamma^j\}=2\eta^{ij},\qquad
\Gamma^{0\cdots10}=-\mathbf1_{32}.
\end{equation}
For an increasing index set, the antisymmetrized gamma matrix equals the ordered product. Accordingly, the slash of a four-form with independent increasing coefficients is simply the sum $\sum_{i<j<k<l}F_{ijkl}\Gamma^{ijkl}$; no extra factorial is inserted.

The antilinear map
\begin{equation}
\mathfrak r(\epsilon)=B_M\epsilon^*,\qquad
B_M=\Gamma^{2468\,10},\qquad
B_MB_M^*=1,\qquad B_M(\Gamma^i)^*=\Gamma^iB_M
\end{equation}
implements the Majorana real structure. All frame, flux and projector coefficients used here are real, so the superconnection and the selected subbundles commute with $\mathfrak r$. A complex solution space of dimension $n$ is therefore the complexification of a real fixed subspace of dimension $n$. The counts sixteen, eight and four in the text are real supercharge counts, not complex counts to be doubled. These local statements do not determine the spin structure or holonomy of an arbitrary quotient.

\bibliographystyle{unsrtnat}
\bibliography{references}
\end{document}
